\documentclass[11pt,letterpaper]{article}
\usepackage[margin=0.7in]{geometry}
\usepackage{times}
\usepackage{titlesec}
\usepackage{enumitem}
\usepackage{hyperref}
\usepackage{setspace}
\usepackage{fancyhdr}
\usepackage{titling}
\usepackage{graphicx}

\setlength{\parindent}{0pt}
\setlength{\parskip}{0.5em}
\onehalfspacing

\titleformat{\section}{\large\bfseries}{\thesection}{0.5em}{}
\titleformat{\subsection}{\normalsize\bfseries}{\thesubsection}{0.5em}{}

\pagestyle{fancy}
\fancyhf{}
\rhead{\thepage}
\lhead{\textit{Aegis — Credential Defense for LLM Agents}}

\begin{document}

\begin{center}
\Large\textbf{Aegis — A Hybrid ML + LLM Credential Defense Layer for LLM Agents}\\[0.3em]
\normalsize\textit{Gauntlet AI Capstone Proposal}
\end{center}

\section*{The Problem}

LLM agents must keep real credentials (API keys, database passwords, OAuth tokens) in the same context window as untrusted content to function. This creates a structural attack surface: indirect prompt injection can cause credential exfiltration. Text-level defenses fail against encodings, slow multi-turn leakage, and — most critically — tool-call arguments. Recent benchmarks show frontier agents achieve only 17\% verified success on related tasks, with up to 40\% of functionally passing patches rejected for cryptographic or semantic violations. The problem is architectural, not incremental.

\section*{Project Direction}

This is a \textbf{Combine ML with LLM applications} project (Direction A). Classical machine learning provides fast, cheap prediction of exfiltration risk from activation features, provenance signals, and behavioral traces. LLM agents handle proposal, repair, and runtime decision-making. Neither paradigm alone solves the problem.

\section*{Technical Approach}

Aegis is a runtime security gateway inspired by Headroom's pipeline model. Messages flow through three stages that treat AI as a first-class primitive:

\textbf{Inspect} — Activation probes, behavioral signals, and tool-argument parsing.  
\textbf{Score} — Classical ML risk model (CIFT-ML) combining provenance and canary signals.  
\textbf{Enforce} — Policy engine that decides Allow / Warn / Block / Escalate.

Aegis supports multiple entry points for easy adoption: SDK wrapper, standalone proxy (zero code changes), and framework plugins. The ML components improve over time from agent trajectories, creating a closed feedback loop. Tool-call argument scanning and cumulative leakage accounting are first-class concerns.

\section*{Scope (Two Weeks)}

\textbf{Minimum shippable}: A working sidecar that intercepts context and tool calls, runs the three layers, and demonstrates measurable improvement over text-only baselines on held-out encoded + multi-turn cases.

\textbf{Ambitious version}: Closed-loop retraining of the ML probe from agent outcomes + first quantitative results on tool-call argument exfiltration.

\section*{Why It's Compelling}

This directly attacks a documented, high-stakes gap that no existing product or research system fully closes. It is visibly ambitious (hybrid closed-loop system on a structural problem), measurable, and timely. Team members will own distinct, high-visibility layers with clear ownership and strong demo potential.

\section*{Hard Parts / Risks}

\begin{itemize}[leftmargin=1.2em]
    \item Tool-call argument scanning has almost no published precedent.
    \item Pre-output detection requires white-box access or strong black-box proxies.
    \item False-positive rate must stay low enough for real use; over-blocking kills utility.
    \item The ML components must demonstrably improve the LLM pipeline, not just add overhead.
\end{itemize}

\section*{Demo Sketch (10 minutes)}

Live side-by-side comparison: baseline agent vs.\ Aegis-protected agent under three attacks (encoded single-turn, low-rate multi-turn dripping, tool-call argument injection). Real-time visualization of the three layers firing, with quantitative results table (detection rate, false blocks, latency) shown live. End with a short architecture diagram and the closed-loop retraining story.

\end{document}